\documentclass[12pt,a4paper]{article}

\usepackage[italian]{babel}
\usepackage{geometry}
\geometry{margin=2.5cm}
\usepackage{array}
\usepackage{tabularx}
\usepackage{xcolor}

\definecolor{negativeRed}{RGB}{180, 40, 40}
\definecolor{orangino}{RGB}{205, 127, 50}

\newcommand{\cognitivetable}[4]{%
	\begin{center}
		\renewcommand{\arraystretch}{1.5}
		\begin{tabularx}{\linewidth}{|X|X|}
			\hline 
			\textbf{Domanda} & \textbf{Osservazioni}
			\\ \hline
			\textbf{1. L'utente sa cosa deve ottenere con l'azione?} (sa cosa deve fare?) & #1 \\ \hline
			\textbf{2. L'utente può individuare lo strumento di interazione?} & #2 \\ \hline
			\textbf{3. L'utente può anche capire che quello è lo strumento giusto per fare ciò che vuol fare?} & #3 \\ \hline
			\textbf{4. Dopo che l'azione è eseguita, l'utente capisce la risposta che ottiene?} & #4 \\ \hline
		\end{tabularx}
	\end{center}
}

\begin{document}
	\section*{Documento di valutazione}
	
	\noindent
	\textbf{Codice di valutazione:} 2 \\
	\textbf{Valutatore:} Lublanis\\
	\textbf{Data di valutazione:} 03/08/2026 \\
	\textbf{Sistema da valutare:} Sito web Home Banking \\
	\textbf{Scopo:} Analizzare le difficoltà dell'utente nell'utilizzare il sistema.\\
	\textbf{Compito:} Reimpostare la password\\
	
	\subsection*{Valutazione Azione}
	
	\textbf{Andare nella sezione del profilo utente}
	\cognitivetable{%
		{Sì. L'utente sa \textcolor{negativeRed}{(forse)} che deve accedere alla sezione dedicata ai suoi dati per trovare le funzionalità relativa alla modifica della password.}%
	}{%
		Sì, il pulsante "Il Mio Profilo" è evidente, è uno dei primi elementi che compaiono a video quando accedo alla dashboard.
	}{%
		Sì, un pulsante giallo "Il Mio Profilo" è abbastanza autoesplicativo per poter comprendere che rimanda ai dettagli del proprio profilo.
	}{%
		Sì. Il sistema mostra la pagina relativa alle informazioni dell'utente, con le relative funzionalità associate.%
	}
	
		\textbf{Andare nella sezione modifica password}
	\cognitivetable{%
		{Sì. Cliccando su modifica password sa che andrà a modificare la password.}%
	}{%
		Sì, il pulsante è molto visibile.
	}{%
		\textcolor{negativeRed}{No, il pulsante è ripetuto due volte e confondibile con un altro pulsante chiamato "Sicurezza".}
	}{%
		Sì. Il sistema mostra la pagina relativa alla modifica della password, con il form dedicato per l'inserimento dei dati (vecchia password, nuova password, conferma nuova password).%
	}
	
	\textbf{Inserire la password}
	\cognitivetable{%
		Sì. L'utente sa che sta inserendo i dati (vecchia password, nuova password, conferma nuova password) sta andando a cambiare la sua password (titolo della sezione in alto).
	}{%
		Sì, il form è al centro della pagina.
	}{%
		Sì, il titolo del form è "Cambio password".
	}{%
		Sì, il sistema comunica eventuali errori o conferma dell'avvenuto cambiamento all'utente.%
	}
	
	\textbf{Confermare l'operazione}
	\cognitivetable{%
		\textcolor{negativeRed}{No, il sistema non chiede conferma all'utente.}
	}{%
		\textcolor{negativeRed}{No, il sistema non chiede conferma all'utente.}
	}{%
		\textcolor{negativeRed}{No, il sistema non chiede conferma all'utente.}
	}{%
		\textcolor{negativeRed}{No, il sistema non chiede conferma all'utente.}
	}
	
\end{document}